\documentclass[
aps,
prl,
groupedaddress,
superscriptaddress,
floatfix,
%showpacs,
%showkeys,
%draft,
%preprint,
%reprint,
%twocolumn,
notitlepage
]{revtex4-1}


\usepackage{graphicx}% Include figure files
%\usepackage[dvips]{graphicx}
\usepackage{amsmath,amssymb}
\usepackage{braket}
%\usepackage{pscyr}
%\usepackage[utf8]{inputenc} % needed for Cyrillic fonts
%\usepackage[T2A]{fontenc} % needed for Cyrillic fonts
\usepackage{hyperref} % hypertext capabilities
\usepackage{subfigure}
%\usepackage{dcolumn} % Align table columns on decimal point
\usepackage{bm} % bold math
\usepackage{mathtools} % dcases (uncompressed fractions)
\usepackage{float}

\usepackage{placeins}


\usepackage[utf8]{inputenc} % needed for Cyrillic fonts
\usepackage[T1,T2A]{fontenc} % needed for Cyrillic fonts
\usepackage[english,russian]{babel} 

\usepackage{float}

% My definitions
\def\sech{{\rm sech}}
\def\rot{{\rm rot}}
\def\div{{\rm div}}
\def\arcsinh{{\rm arcsinh}}
\def\Re{{\rm Re\,}}
\def\Im{{\rm Im\,}}
\def\arccot{{\rm arccot}}
\newcommand{\p}{\partial}
\newcommand{\ep}{\varepsilon}
\newcommand{\br}{\break}
\newcommand{\om}{\omega}
\newcommand{\ph}{\varphi}
\newcommand{\nn}{\nonumber}
\newcommand{\ka}{\kappa}
\newcommand{\al}{\alpha}
\newcommand{\la}{\lambda}
\newcommand{\be}{\beta}

\begin{document}

\title{Исследование нелинейных топологических состояний вещества с помощью экситон-поляритонов: краевые солитоны в решетке кагоме}

\author{D. R. Gulevich} 
\affiliation{ITMO University, St. Petersburg 197101, Russia}
\affiliation{Department of Physics, University of Bath, Bath BA2 7AY, United Kingdom}

\author{D. Yudin} 
\affiliation{ITMO University, St. Petersburg 197101, Russia}
\affiliation{Division of Physics and Applied Physics, Nanyang Technological University 637371, Singapore}

\author{D. V. Skryabin} 
\affiliation{ITMO University, St. Petersburg 197101, Russia}
\affiliation{Department of Physics, University of Bath, Bath BA2 7AY, United Kingdom}

\author{I. V. Iorsh} 
\affiliation{ITMO University, St. Petersburg 197101, Russia}
\affiliation{Division of Physics and Applied Physics, Nanyang Technological University 637371, Singapore}

\author{I. A. Shelykh} 
\affiliation{ITMO University, St. Petersburg 197101, Russia}
\affiliation{Science Institute, University of Iceland, Dunhagi 3, IS-107, Reykjavik, Iceland}

%\keywords{Keyword1, Keyword2, Keyword3}

\begin{abstract}
Вещество в нетривиальной топологической фазе обладает уникальными свойствами, такими как поддержка однонаправленных краевых мод на границе раздела. Именно существование таких мод отвечает за замечательные свойства топологического изолятора - материала, который изолирует в объеме, но проводит на его поверхности, наряду со многими из его недавно предложенных фотонных и поляритонных аналогов. Показано, что экситон-поляритонная жидкость в нетривиальной топологической фазе в решетке кагоме поддерживает нелинейные возбуждения в виде солитонов, построенных из волновых пакетов краевых топологических мод - краевых топологических солитонов. Наши теоретические и численные результаты указывают на появление ярких, темных и серых солитонов, находящихся вблизи границы полосы решетки. В параболической области дисперсии солитоны описываются огибающими функциями, удовлетворяющими нелинейному уравнению Schr\"odinger 
При столкновении несколько топологических краевых солитонов возникают неискаженными, что доказывает, что они являются истинными солитонами в отличие от уединенных волн, для которых такое требование не требуется. Важно отметить, что решетка кагоме поддерживает топологический краевой режим с нулевой групповой скоростью в отличие от других типов усеченных решеток. Это дает более тонкий контроль над скоростью солитона, которая может принимать как положительные, так и отрицательные значения в зависимости от выбора формирования его топологических краевых мод.
\end{abstract}


%\flushbottom
%\maketitle
%\thispagestyle{empty}

%\maketitle
{\let\newpage\relax\maketitle}

%-----------------------------------------------------------------------------------------------
%-----------------------------------------------------------------------------------------------
\section*{Введение}
%-----------------------------------------------------------------------------------------------
%-----------------------------------------------------------------------------------------------

Различные физические системы часто демонстрируют сходство в лежащих в их основе физических явлениях, что приводит к идее использования хорошо управляемых систем для имитации свойств менее контролируемых и менее доступных. .
Ярким примером является сходство электронных и фотонных систем~\cite{Georgescu2014}.
Поэтому неудивительно, что с быстрым развитием топологических изоляторов в контексте электронных систем~\cite{Kane-Mele,Bernevig,Konig}топологические идеи были
также широко исследуется в фотонных системах. Среди новаторских работ -
исследование киральных краевых состояний в фотонных кристаллах~\cite{Raghu-2008} где кривизна Берри для фотонных зон была введена по аналогии с электронными системами, фотонными аналогами эффекта Холла~\cite{Haldane2008a, Wang2008,Wang-2009} и топологическими изоляторами~\cite{Hafezi-2011,Fang-2012,Rechtsman-2013,Khanikaev2013}, а также целым рядом как теоретических и экспериментальные работы, демонстрирующие эффекты нетривиальной топологии в электромагнитных системах в диапазоне частот от радио до оптики~\cite{Lu-2014}.

Однако, несмотря на этот успех, оптические схемы плохо подходят для непосредственного включения нелинейных эффектов, в то время как реализация нарушения симметрии относительно обращения времени, общей составляющей топологических фаз, остается сложной задачей.
В этом отношении системы на основе экситон-поляритонов~\cite{Carusotto2013}, квазичастиц, возникающих в результате сильной связи экситонов квантовой ямы и фотонов резонатора в микрорезонаторах, являются преимуществом.
Будучи гибридными возбуждениями легкого вещества, их фотонные свойства позволяют эффективно управлять с использованием профиля оптического потенциала, в то время как их экситонная природа приводит к значительным взаимодействиям и сильному нелинейному отклику~\cite{CedraMendez2010,Kim2013,Jacqmin2014,Baboux2016}.
Более того, из-за спина экситона симметрия экситон-поляритонной системы относительно обращения времени может быть легко нарушена приложением внешнего магнитного поля. Это делает поляритонные системы привлекательными как с точки зрения приложений в перспективных поляритонных устройствах~\cite{Liew2011,Sanvitto2016} и в качестве уникальной лаборатории для моделирования топологических свойств вещества.

Было высказано несколько предложений по созданию топологически нетривиальных состояний экситон-поляритонов. В последние годы появлению нетривиальной топологии и существованию топологически защищенных краевых состояний в поляритонных решетках различной геометрии посвящен ряд работ~\cite{Karzig-PRX-2015, Bardyn-PRB-2015, Nalitov-Z, Yi-PRB-2016, Bardyn-PRB-2016, Janot-PRB-2016, Gulevich-kagome}.
%It is, therefore, natural, that the focus of attention in the study of the effects of non-trivial topology shifts towards systems with nonlinearity, where exciton-polaritons play an important role.
В настоящее время фокус внимания при изучении эффектов нетривиальной топологии смещается в сторону систем с нелинейностью, в которых экситон-поляритоны, благодаря своим уникальным свойствам, играют особую роль. Среди недавних работ, в которых исследовалось взаимодействие нелинейных эффектов с нетривиальной топологией, можно назвать исследования автолокализованных состояний~\cite{Lumer-PRL-2013, Ostrovskaya2013}, самоиндуцированных топологических переходов~\cite{Hadad-PRB-2016}, топологических боголюбовских возбуждений~\cite{Bardyn-PRB-2016},
подавление топологических фаз~\cite{Bleu-PRB-2016}, вихрей в решетках~\cite{Kartashov-OL-2016}, спин-мейсснеровских состояний в кольцевых резонаторах~\cite{DGulevich-Meissner}, солитонов в решетках~\cite{Ablowitz-PRA-2014, Leykam-PRL-2016, Kartashov-Optica-2016, CedraMendez2016} и димеров
цепи~\cite{Soln-PRL-2017}.
В исх.~\cite{Gulevich-kagome} было замечено, что в определенном диапазоне параметров решетка кагоме обладает сильно нелинейной дисперсией топологического краевого состояния с хорошо выраженным минимумом и максимумом внутри объемной щели. В настоящей работе мы показываем, что такая своеобразная дисперсия приводит к появлению нелинейных краевых возбуждений в виде солитонов. Такие возбуждения оказываются истинными солитонами в отличие от уединенных волн, для которых не требуется восстановление формы при столкновении~\cite{Rajaraman}. В отличие от солитонов в сотовой решетке~\cite{Kartashov-Optica-2016} 
скорость топологических солитонов краевого состояния в решетке кагоме может принимать значения в широком диапазоне от положительных до отрицательных.
в зависимости от выбора квазиимпульсов составляющих топологических краевых мод.
%be tuned from positive to negative by controlling quasimomenta of the constituting topological edge modes.
%central quasimomentum of wavepacket.

%-----------------------------------------------------------------------------------------------
\subsection*{Теория нелинейных топологических краевых возбуждений в двумерной решетке}
%-----------------------------------------------------------------------------------------------

Для изучения нелинейных топологических краевых возбуждений мы используем приближение среднего поля. Мы рассматриваем полосу решетки кагоме с $M$ в ее элементарной ячейке, см. Рис.~\ref{fig:strip-disp}a, и вводим векторы $\pmb{\psi}_m^{\sigma}=(\psi_{m,1}^{\sigma},...,\psi_{m,M}^{\sigma})^T$ составленные из спинорных компонентов $\sigma=\pm$ в каждый узел $m$-й элементарной ячейки. , где индекс $m$ нумерует элементарные ячейки в направлении $x$ (см. рис.~\ref{fig:strip-disp}a).
Тогда временная эволюция $m$-го вектора описывается выражением
связанная система уравнений
\begin{equation}
i\p_t\pmb{\psi}_m^{\sigma}=\sum_{\tau}\left(\hat H_{-1}^{\sigma,\tau}\pmb{\psi}_{m-1}^{\tau}+\hat H_0^{\sigma,\tau}\pmb{\psi}_m^{\tau}+\hat H_{+1}^{\sigma,\tau}\pmb{\psi}_{m+1}^{\tau}\right)+ \hat{N}^{\sigma}(\pmb{\psi}_m^{+},\pmb{\psi}_m^{-}) \pmb{\psi}_m^{\sigma},
\label{td-problem}
\end{equation} 
где $\hat{N}^{\sigma}(\pmb{\psi}_m^{+},\pmb{\psi}_m^{-})$ - диагональная матрица с элементами
$$\left[\hat{N}^{\sigma}(\pmb{\psi}_m^{+},\pmb{\psi}_m^{-})\right]_{ij} = \delta_{i,j}\left(|\psi_{m,j}^{\sigma}|^2 + \alpha |\psi_{m,j}^{\bar\sigma}|^2\right),$$
параметр $\alpha\equiv \alpha_2/\alpha_1$, а 
$\pmb{\psi}_m^{\sigma}$ нормированы на $\sum_{\sigma,m}\langle\pmb{\psi}_m^{\sigma},\pmb{\psi}_m^{\sigma}\rangle \equiv \sum_{\sigma,m,j}[\pmb{\psi}_m^{\sigma*}]_j [\pmb{\psi}_m^{\sigma}]_j=N/\alpha_1$ 
% N J/\alpha_1
где $N$ - полное количество поляритонов в системе.
Матрицы $\hat H_{j}^{\sigma,\tau}$удовлетворяют $(\hat H_{j}^{\sigma,\tau})^\dagger=\hat H_{-j}^{\tau,\sigma}$
для всех $\sigma,\tau=\pm$, $j=0,\pm1$, что обеспечивает эрмитичность гамильтониана.
В отсутствие взаимодействий стационарные состояния~\eqref{td-problem} являются блоховскими волнами $e^{-i\mu t+ikLm}\,\pmb{u}_{k}^{\sigma}$ где $L$ - пространственная протяженность элементарной ячейки в направлении $x$ а $\pmb{u}_{k}^{\sigma}$ - %$m$-independent 
решения проблемы собственных значений
\begin{equation}
\sum_{\tau}\left(e^{-ikL} \hat H_{-1}^{\sigma,\tau}+\hat H_0^{\sigma,\tau}+e^{ikL} \hat H_{+1}^{\sigma,\tau}\right) \pmb{u}_{k,n}^{\tau}=\mu_{k,n}\,\pmb{u}_{k,n}^{\sigma}
\quad\text{для $\sigma=\pm$,}
\label{operator}
\end{equation}
где $n$ перечисляет энергетические зоны. Поскольку оператор в левой части~\eqref{operator} является самосопряженным, $\pmb{u}_{k,n}^{\sigma}$ образует ортонормированное множество, 
$$
\sum_\sigma\langle\pmb{u}_{k,n}^{\sigma},\pmb{u}_{k,m}^{\sigma}\rangle 
=\sum_{\sigma,j}[\pmb{u}_{k,n}^{\sigma*}]_j[\pmb{u}_{k,m}^{\sigma}]_j=\delta_{n,m}.
$$
Пусть $n=n_e$ - индекс зоны, соответствующий одной из дисперсий ПЭМ, как на рис.~\ref{fig:strip-disp}б.
Затем решение полной проблемы~\eqref{td-problem} можно искать в виде волнового пакета, сосредоточенного вокруг $k_e$,
\begin{equation}
\pmb{\psi}^{\sigma}_m(k_e,t)
=\sum_n \int_{-\pi/L}^{\pi/L} A_n(\kappa,t)\pmb{u}^{\sigma}_{k_e+\kappa,n} e^{i(k_e+\kappa)Lm} d\kappa
\approx \int_{-\pi/L}^{\pi/L} A(\kappa,t)\pmb{u}^{\sigma}_{k_e+\kappa,n_e} e^{i(k_e+\kappa)Lm} d\kappa,
\label{psi-anzatz}
\end{equation}
где мы предположили, что один ПЭМ с $n=n_e$ и амплитудой $A(\kappa,t)$ доминирует над другими модами.
Для $A(\kappa,t)$ 
необходимо зафиксировать калибровку базисных векторов $\pmb{u}^{\sigma}_{k,n_e}$. Всегда можно выбрать такую калибровку, чтобы уравнение
\begin{equation}
\sum_{\sigma}\langle\pmb{u}^{\sigma}_{k,n_e},\frac{\p }{\p k}\pmb{u}^{\sigma}_{k,n_e}\rangle=0
\label{gauge}
\end{equation}
выполнялось, по крайней мере, в пределах 
%arbitrary 
открытого интервала $k$ меньшего, чем зона Бриллюэна. В то время как действительная часть~\eqref{gauge} гарантируется нормировкой, мнимая часть~\eqref{gauge} может быть $U(1)$ фаз собственных векторов.
Подставляя~\eqref{psi-anzatz} в~\eqref{td-problem} и формируя скалярное произведение с $\pmb{u}^{\sigma}_{k_e,n_e}$ мы получаем
\begin{equation}
\int_{-\pi/L}^{\pi/L}\sum_\sigma
\left[ \left(i\frac{\p A(\kappa,t)}{\p t}-\mu_{k_e+\kappa,n_e}A(\kappa,t)\right)\langle\pmb{u}^{\sigma}_{k_e,n_e},\pmb{u}^{\sigma}_{k_e+\kappa,n_e}\rangle
- A(\kappa,t)\langle\pmb{u}^{\sigma}_{k_e,n_e},\hat{N}^{\sigma}(\pmb{\psi}_m^{+},\pmb{\psi}_m^{-})\pmb{u}^{\sigma}_{k_e+\kappa,n_e}\rangle 
\right] e^{i\kappa Lm} d\kappa = 0
\label{nls-0}
\end{equation}
для всех $m$.
Дифференцируя~\eqref{gauge} получаем
\begin{equation}
\sum_{\sigma}\langle\frac{\p }{\p k}\pmb{u}^{\sigma}_{k,n_e},\frac{\p }{\p k}\pmb{u}^{\sigma}_{k,n_e}\rangle+
\sum_{\sigma}\langle\pmb{u}^{\sigma}_{k,n_e},\frac{\p^2 }{\p^2 k}\pmb{u}^{\sigma}_{k,n_e}\rangle=0.
\label{p2k}
\end{equation}
. Калибровочное условие~\eqref{gauge} 
%together with the completeness of the basis vectors 
означает, что производная $\p\pmb{u}^{\sigma}_{k,n_e}/\p k$ включает возбуждения всех других зон, кроме $n_e$ (в общем, объемных зон и других ветвей ПЭМ). В частности, в $k=k_e$,
\begin{equation}
\frac{\p }{\p k}\pmb{u}^{\sigma}_{k,n_e}\Big|_{k=k_e} = \sum_{n\neq n_e} c_n \pmb{u}^{\sigma}_{k_e,n}.
\label{cn}
\end{equation}
Разлагая $\pmb{u}^{\sigma}_{k_e+\kappa,n_e}$ в ряд Маклорена в $\kappa$
и используя~\eqref{p2k} и~\eqref{cn}, мы получаем
\begin{equation}
\sum_{\sigma}\langle\pmb{u}^{\sigma}_{k_e,n_e},\pmb{u}^{\sigma}_{k_e+\kappa,n_e}\rangle\approx 1-\frac12 \kappa^2\sum_{n\neq n_e}|c_n|^2.
\label{cond}
\end{equation}
Предполагая, что возбуждение полос $n\neq n_e$ слабое, что гарантируется большим энергетическим отрывом полосы $n_e$ от остальных полос, можно пренебречь
вклад объемных зон в~\eqref{nls-0}.
Разлагая $\mu_{k_e+\kappa,n_e}$ на ряд Маклорена до 2-го порядка в $\kappa$ и интегрируя, уравнение~\eqref{nls-0} сокращается до
\begin{equation}
i\frac{\p \tilde{A}}{\p t} =
\sum_{n=0}^{\infty} \frac{(-i)^n}{n!} \mu_{k_e}^{(n)}\frac{\partial^n \tilde{A}}{\partial x^n}\Big|_{x=Lm}
+ g |\tilde{A}|^2 \tilde{A}\quad\text{для всех $x=Lm$},
\label{nls-2}
\end{equation}
где 
$$
\tilde{A}(x,t)\equiv \int_{-\pi/L}^{\pi/L} A(\kappa,t) e^{i\kappa x}  d\kappa
$$
является функцией времени и непрерывной переменной $x$, а 
$$
g\equiv \sum_\sigma\langle\pmb{u}^{\sigma}_{k_e,n_e},\hat{N}^{\sigma}(\pmb{u}_{k_e,n_e}^+,\pmb{u}_{k_e,n_e}^{-})\pmb{u}^{\sigma}_{k_e,n_e}\rangle
$$
является эффективным параметром нелинейности. Для $\alpha \ge -1$, которое является обычным случаем для поляритонных взаимодействий~\cite{Vladimirova}, можно доказать, что $g$ принимает только неотрицательные значения, $g\ge 0$, независимо от~$\pmb{u}_{k_e,n_e}^\sigma$, тем самым описывая дефокусирующую нелинейность.
Наконец, 
\begin{equation}
\pmb{\psi}^{\sigma}_m(k_e,t) \approx  \tilde{A}(Lm,t)\, e^{i k_e Lm}\, \pmb{u}^{\sigma}_{k_e,n_e}.
\label{psi-final}
\end{equation}

Используя~\eqref{cond}, мы можем проанализировать справедливость примененного приближения.
Коэффициенты $c_n$ могут быть вычислены из стандартной теории возмущений~\cite{LL-III}, что дает
\begin{equation}
\quad
c_n = \frac{1}{\mu_{k_e,n_e}-\mu_{k_e,n}}\,
\sum_{\sigma,\tau} \langle \pmb{u}^{\sigma}_{k_e,n}, \hat{V}^{\sigma,\tau}
\pmb{u}^{\tau}_{k_e,n_e}
\rangle,
\quad
\hat{V}^{\sigma,\tau} = iL \left(e^{ik_eL} \hat H_{+1}^{\sigma,\tau} - e^{-ik_eL} \hat H_{-1}^{\sigma,\tau}\right).
\label{cn-pert}
\end{equation}
Таким образом, мы можем оценить
\begin{multline}
\sum_{n\neq n_e}|c_n|^2 \le \frac{1}{\Delta\mu^2} 
\sum_{n\neq n_e} 
\left|\sum_{\sigma,\tau} \langle \pmb{u}^{\sigma}_{k_e,n}, \hat{V}^{\sigma,\tau}\pmb{u}^{\tau}_{k_e,n_e} \rangle\right|^2
=
\frac{1}{\Delta\mu^2} \left(
\sum_{\sigma,\tau,\tau'} 
\langle \hat{V}^{\sigma,\tau}\pmb{u}^{\tau}_{k_e,n_e} , \hat{V}^{\sigma,\tau'}\pmb{u}^{\tau'}_{k_e,n_e}\rangle
- \left| \sum_{\sigma,\tau} \langle \pmb{u}^{\sigma}_{k_e,n_e}, \hat{V}^{\sigma,\tau}\pmb{u}^{\tau}_{k_e,n_e} \rangle \right|^2
\right) 
\\
\le \frac{1}{\Delta\mu^2}
\sum_{\sigma,\tau,\tau'} 
\langle \hat{V}^{\sigma,\tau}\pmb{u}^{\tau}_{k_e,n_e} , \hat{V}^{\sigma,\tau'}\pmb{u}^{\tau'}_{k_e,n_e}\rangle,
\label{estimate}
\end{multline}
где мы ввели разделение по энергии от объемных мод $\Delta\mu\equiv {\rm min}_{n\neq n_e} |\mu-\mu_n|$
и использовали полноту базисного набора.
Последний член в~\eqref{estimate} можно оценить, используя точную форму~\eqref{cn-pert} для оператора $\hat{V}^{\sigma,\tau}$ сохраняя только прыжковые члены с сохраняющейся круговой поляризацией в исходном гамильтониане.
Это дает оценку $\sim L^2/\Delta\mu^2$.
Определяя среднее значение
$$
\langle f(\kappa) \rangle_{A} \equiv \left( \int_{-\pi/L}^{\pi/L} |A(\kappa,t)| d\kappa\right)^{-1} \int_{-\pi/L}^{\pi/L} |A(\kappa,t)| f(\kappa) d\kappa
$$
и используя~\eqref{cond},~\eqref{estimate} и $L=2$ мы можем получить критерии применимости~\eqref{nls-2} и~\eqref{psi-final},
\begin{equation}
2\langle \kappa^2 \rangle_{A} \ll \Delta\mu^2,
\label{cond-final}
\end{equation}
которые являются ограничением на протяженность волнового пакета в k- пространство для заданного энергетического отделения от объемных мод.
%size of the energy gap.
Отметим, что очевидное требование независимости~\eqref{psi-final} от выбора элементарной ячейки полоски кагоме (заметим, что выбор элементарной ячейки на рис.~\ref{fig:strip-disp} не уникален) требует либо слабой зависимости из $\tilde{A}(x,t)$ по $x$ (то есть небольшой $\langle \kappa^2 \rangle_{A}$), или сильная локализация краевой моды
вблизи границы в пределах первых нескольких рядов (обеспечивается большой $\Delta\mu^2$).

%-----------------------------------------------------------------------------------------------
%-----------------------------------------------------------------------------------------------
%\bibliography{libkagome}
%-----------------------------------------------------------------------------------------------
%-----------------------------------------------------------------------------------------------
%\noindent

\begin{thebibliography}{10}
	\expandafter\ifx\csname url\endcsname\relax
	\def\url#1{\texttt{#1}}\fi
	\expandafter\ifx\csname urlprefix\endcsname\relax\def\urlprefix{URL }\fi
	\expandafter\ifx\csname doiprefix\endcsname\relax\def\doiprefix{DOI }\fi
	\providecommand{\bibinfo}[2]{#2}
	\providecommand{\eprint}[2][]{\url{#2}}
	
	\bibitem{Georgescu2014}
	\bibinfo{author}{Georgescu, I.}, \bibinfo{author}{Ashhab, S.} \&
	\bibinfo{author}{Nori, F.}
	\newblock \bibinfo{title}{Quantum simulation}.
	\newblock \emph{\bibinfo{journal}{Reviews of Modern Physics}}
	\textbf{\bibinfo{volume}{86}}, \bibinfo{pages}{153} (\bibinfo{year}{2014}).
	
	\bibitem{Kane-Mele}
	\bibinfo{author}{Kane, C.~L.} \& \bibinfo{author}{Mele, E.~J.}
	\newblock \bibinfo{title}{Quantum spin hall effect in graphene}.
	\newblock \emph{\bibinfo{journal}{Phys. Rev. Lett.}}
	\textbf{\bibinfo{volume}{95}}, \bibinfo{pages}{226801}
	(\bibinfo{year}{2005}).
	
	\bibitem{Bernevig}
	\bibinfo{author}{Bernevig, B.~A.}, \bibinfo{author}{Hughes, T.~L.} \&
	\bibinfo{author}{Zhang, S.~C.}
	\newblock \bibinfo{title}{Quantum spin hall effect and topological phase
		transition in hgte quantum wells}.
	\newblock \emph{\bibinfo{journal}{Science}} \textbf{\bibinfo{volume}{314}},
	\bibinfo{pages}{1757} (\bibinfo{year}{2006}).
	
	\bibitem{Konig}
	\bibinfo{author}{K\"{o}nig, M.} \emph{et~al.}
	\newblock \bibinfo{title}{Quantum spin hall insulator state in hgte quantum
		wells}.
	\newblock \emph{\bibinfo{journal}{Science}} \textbf{\bibinfo{volume}{318}},
	\bibinfo{pages}{766} (\bibinfo{year}{2007}).
	
	\bibitem{Raghu-2008}
	\bibinfo{author}{Raghu, S.} \& \bibinfo{author}{Haldane, F. D.~M.}
	\newblock \bibinfo{title}{Analogs of quantum-hall-effect edge states in
		photonic crystals}.
	\newblock \emph{\bibinfo{journal}{Phys. Rev. A}} \textbf{\bibinfo{volume}{78}},
	\bibinfo{pages}{033834} (\bibinfo{year}{2008}).
	
	\bibitem{Haldane2008a}
	\bibinfo{author}{Haldane, F.} \& \bibinfo{author}{Raghu, S.}
	\newblock \bibinfo{title}{Possible realization of directional optical
		waveguides in photonic crystals with broken time-reversal symmetry}.
	\newblock \emph{\bibinfo{journal}{Phys. Rev. Lett.}}
	\textbf{\bibinfo{volume}{100}}, \bibinfo{pages}{013904}
	(\bibinfo{year}{2008}).
	
	\bibitem{Wang2008}
	\bibinfo{author}{Wang, Z.}, \bibinfo{author}{Chong, Y.},
	\bibinfo{author}{Joannopoulos, J.~D.} \& \bibinfo{author}{Solja{\v{c}}i{\'c},
		M.}
	\newblock \bibinfo{title}{Reflection-free one-way edge modes in a gyromagnetic
		photonic crystal}.
	\newblock \emph{\bibinfo{journal}{Phys. Rev. Lett.}}
	\textbf{\bibinfo{volume}{100}}, \bibinfo{pages}{013905}
	(\bibinfo{year}{2008}).
	
	\bibitem{Wang-2009}
	\bibinfo{author}{Wang, Z.}, \bibinfo{author}{Chong, Y.},
	\bibinfo{author}{Joannopoulos, F.~D.} \& \bibinfo{author}{Solja\v{c}i\'{c},
		M.}
	\newblock \bibinfo{title}{Observation of unidirectional backscattering-immune
		topological electromagnetic states}.
	\newblock \emph{\bibinfo{journal}{Nature}} \textbf{\bibinfo{volume}{461}},
	\bibinfo{pages}{772--775} (\bibinfo{year}{2009}).
	
	\bibitem{Hafezi-2011}
	\bibinfo{author}{Hafezi, M.}, \bibinfo{author}{Demler, E.~A.},
	\bibinfo{author}{Lukin, M.~D.} \& \bibinfo{author}{Taylor, J.~M.}
	\newblock \bibinfo{title}{Robust optical delay lines with topological
		protection}.
	\newblock \emph{\bibinfo{journal}{Nature Physics}}
	\textbf{\bibinfo{volume}{7}}, \bibinfo{pages}{907--912}
	(\bibinfo{year}{2011}).
	
	\bibitem{Fang-2012}
	\bibinfo{author}{Fang, K.}, \bibinfo{author}{Yu, Z.} \& \bibinfo{author}{Fan,
		S.}
	\newblock \bibinfo{title}{Realizing effective magnetic field for photons by
		controlling the phase of dynamic modulation}.
	\newblock \emph{\bibinfo{journal}{Nature Photonics}}
	\textbf{\bibinfo{volume}{6}}, \bibinfo{pages}{782} (\bibinfo{year}{2012}).
	
	\bibitem{Rechtsman-2013}
	\bibinfo{author}{Rechtsman, M.~C.} \emph{et~al.}
	\newblock \bibinfo{title}{Photonic floquet topological insulators}.
	\newblock \emph{\bibinfo{journal}{Nature}} \textbf{\bibinfo{volume}{496}},
	\bibinfo{pages}{196} (\bibinfo{year}{2013}).
	
	\bibitem{Khanikaev2013}
	\bibinfo{author}{Khanikaev, A.~B.} \emph{et~al.}
	\newblock \bibinfo{title}{Photonic topological insulators}.
	\newblock \emph{\bibinfo{journal}{Nature Materials}}
	\textbf{\bibinfo{volume}{12}}, \bibinfo{pages}{233--239}
	(\bibinfo{year}{2013}).
	
	\bibitem{Lu-2014}
	\bibinfo{author}{Lu, L.}, \bibinfo{author}{Joannopoulos, J.~D.} \&
	\bibinfo{author}{Solja\v{c}i\'{c}, M.}
	\newblock \bibinfo{title}{Topological photonics}.
	\newblock \emph{\bibinfo{journal}{Nature Photonics}}
	\textbf{\bibinfo{volume}{8}}, \bibinfo{pages}{821} (\bibinfo{year}{2014}).
	
	\bibitem{Carusotto2013}
	\bibinfo{author}{Carusotto, I.} \& \bibinfo{author}{Ciuti, C.}
	\newblock \bibinfo{title}{Quantum fluids of light}.
	\newblock \emph{\bibinfo{journal}{Reviews of Modern Physics}}
	\textbf{\bibinfo{volume}{85}}, \bibinfo{pages}{299} (\bibinfo{year}{2013}).
	
	\bibitem{CedraMendez2010}
	\bibinfo{author}{Cerda-M\'endez, E.~A.} \emph{et~al.}
	\newblock \bibinfo{title}{Polariton condensation in dynamic acoustic lattices}.
	\newblock \emph{\bibinfo{journal}{Phys. Rev. Lett.}}
	\textbf{\bibinfo{volume}{105}}, \bibinfo{pages}{116402}
	(\bibinfo{year}{2010}).
	
	\bibitem{Kim2013}
	\bibinfo{author}{Kim, N.~Y.} \emph{et~al.}
	\newblock \bibinfo{title}{Exciton–polariton condensates near the dirac point
		in a triangular lattice}.
	\newblock \emph{\bibinfo{journal}{New Journal of Physics}}
	\textbf{\bibinfo{volume}{15}}, \bibinfo{pages}{035032}
	(\bibinfo{year}{2013}).
	
	\bibitem{Jacqmin2014}
	\bibinfo{author}{Jacqmin, T.} \emph{et~al.}
	\newblock \bibinfo{title}{Direct observation of dirac cones and a flatband in a
		honeycomb lattice for polaritons}.
	\newblock \emph{\bibinfo{journal}{Phys. Rev. Lett.}}
	\textbf{\bibinfo{volume}{112}}, \bibinfo{pages}{116402}
	(\bibinfo{year}{2014}).
	
	\bibitem{Baboux2016}
	\bibinfo{author}{Baboux, F.} \emph{et~al.}
	\newblock \bibinfo{title}{Bosonic condensation and disorder-induced
		localization in a flat band}.
	\newblock \emph{\bibinfo{journal}{Phys. Rev. Lett.}}
	\textbf{\bibinfo{volume}{116}}, \bibinfo{pages}{066402}
	(\bibinfo{year}{2016}).
	
	\bibitem{Liew2011}
	\bibinfo{author}{Liew, T.}, \bibinfo{author}{Shelykh, I.} \&
	\bibinfo{author}{Malpuech, G.}
	\newblock \bibinfo{title}{Polaritonic devices}.
	\newblock \emph{\bibinfo{journal}{Physica E}} \textbf{\bibinfo{volume}{43}},
	\bibinfo{pages}{1543} (\bibinfo{year}{2011}).
	
	\bibitem{Sanvitto2016}
	\bibinfo{author}{D.~Sanvitto, D.} \& \bibinfo{author}{K\'{e}na-Cohen, S.}
	\newblock \bibinfo{title}{The road towards polaritonic devices}.
	\newblock \emph{\bibinfo{journal}{Nature Materials}}
	\textbf{\bibinfo{volume}{15}}, \bibinfo{pages}{1061} (\bibinfo{year}{2016}).
	
	\bibitem{Karzig-PRX-2015}
	\bibinfo{author}{Karzig, T.}, \bibinfo{author}{Bardyn, C.-E.},
	\bibinfo{author}{Lindner, N.~H.} \& \bibinfo{author}{Refael, G.}
	\newblock \bibinfo{title}{Topological polaritons}.
	\newblock \emph{\bibinfo{journal}{Phys. Rev. X}} \textbf{\bibinfo{volume}{5}},
	\bibinfo{pages}{031001} (\bibinfo{year}{2015}).
	
	\bibitem{Bardyn-PRB-2015}
	\bibinfo{author}{Bardyn, C.-E.}, \bibinfo{author}{Karzig, T.},
	\bibinfo{author}{Refael, G.} \& \bibinfo{author}{Liew, T. C.~H.}
	\newblock \bibinfo{title}{Topological polaritons and excitons in garden-variety
		systems}.
	\newblock \emph{\bibinfo{journal}{Phys. Rev. B}} \textbf{\bibinfo{volume}{91}},
	\bibinfo{pages}{161413} (\bibinfo{year}{2015}).
	
	\bibitem{Nalitov-Z}
	\bibinfo{author}{Nalitov, A.~V.}, \bibinfo{author}{Solnyshkov, D.~D.} \&
	\bibinfo{author}{Malpuech, G.}
	\newblock \bibinfo{title}{Polariton $\mathbb{Z}$ topological insulator}.
	\newblock \emph{\bibinfo{journal}{Phys. Rev. Lett.}}
	\textbf{\bibinfo{volume}{114}}, \bibinfo{pages}{116401}
	(\bibinfo{year}{2015}).
	
	\bibitem{Yi-PRB-2016}
	\bibinfo{author}{Yi, K.} \& \bibinfo{author}{Karzig, T.}
	\newblock \bibinfo{title}{Topological polaritons from photonic dirac cones
		coupled to excitons in a magnetic field}.
	\newblock \emph{\bibinfo{journal}{Phys. Rev. B}} \textbf{\bibinfo{volume}{93}},
	\bibinfo{pages}{104303} (\bibinfo{year}{2016}).
	
	\bibitem{Bardyn-PRB-2016}
	\bibinfo{author}{Bardyn, C.-E.}, \bibinfo{author}{Karzig, T.},
	\bibinfo{author}{Refael, G.} \& \bibinfo{author}{Liew, T. C.~H.}
	\newblock \bibinfo{title}{Chiral bogoliubov excitations in nonlinear bosonic
		systems}.
	\newblock \emph{\bibinfo{journal}{Phys. Rev. B}} \textbf{\bibinfo{volume}{93}},
	\bibinfo{pages}{020502} (\bibinfo{year}{2016}).
	
	\bibitem{Janot-PRB-2016}
	\bibinfo{author}{Janot, A.}, \bibinfo{author}{Rosenow, B.} \&
	\bibinfo{author}{Refael, G.}
	\newblock \bibinfo{title}{Topological polaritons in a quantum spin hall
		cavity}.
	\newblock \emph{\bibinfo{journal}{Phys. Rev. B}} \textbf{\bibinfo{volume}{93}},
	\bibinfo{pages}{161111} (\bibinfo{year}{2016}).
	
	\bibitem{Gulevich-kagome}
	\bibinfo{author}{Gulevich, D.~R.}, \bibinfo{author}{Yudin, D.},
	\bibinfo{author}{Iorsh, I.~V.} \& \bibinfo{author}{Shelykh, I.~A.}
	\newblock \bibinfo{title}{Kagome lattice from an exciton-polariton
		perspective}.
	\newblock \emph{\bibinfo{journal}{Phys. Rev. B}} \textbf{\bibinfo{volume}{94}},
	\bibinfo{pages}{115437} (\bibinfo{year}{2016}).
	
	\bibitem{Lumer-PRL-2013}
	\bibinfo{author}{Lumer, Y.}, \bibinfo{author}{Plotnik, Y.},
	\bibinfo{author}{Rechtsman, M.~C.} \& \bibinfo{author}{Segev, M.}
	\newblock \bibinfo{title}{Self-localized states in photonic topological
		insulators}.
	\newblock \emph{\bibinfo{journal}{Phys. Rev. Lett.}}
	\textbf{\bibinfo{volume}{111}}, \bibinfo{pages}{243905}
	(\bibinfo{year}{2013}).
	
	\bibitem{Ostrovskaya2013}
	\bibinfo{author}{Ostrovskaya, E.~A.}, \bibinfo{author}{Abdullaev, J.},
	\bibinfo{author}{Fraser, M.~D.}, \bibinfo{author}{Desyatnikov, A.~S.} \&
	\bibinfo{author}{Kivshar, Y.~S.}
	\newblock \bibinfo{title}{Self-localization of polariton condensates in
		periodic potentials}.
	\newblock \emph{\bibinfo{journal}{Phys. Rev. Lett.}}
	\textbf{\bibinfo{volume}{110}}, \bibinfo{pages}{170407}
	(\bibinfo{year}{2013}).
	
	\bibitem{Hadad-PRB-2016}
	\bibinfo{author}{Hadad, Y.}, \bibinfo{author}{Khanikaev, A.~B.} \&
	\bibinfo{author}{Al\`u, A.}
	\newblock \bibinfo{title}{Self-induced topological transitions and edge states
		supported by nonlinear staggered potentials}.
	\newblock \emph{\bibinfo{journal}{Phys. Rev. B}} \textbf{\bibinfo{volume}{93}},
	\bibinfo{pages}{155112} (\bibinfo{year}{2016}).
	
	\bibitem{Bleu-PRB-2016}
	\bibinfo{author}{Bleu, O.}, \bibinfo{author}{Solnyshkov, D.~D.} \&
	\bibinfo{author}{Malpuech, G.}
	\newblock \bibinfo{title}{Interacting quantum fluid in a polariton chern
		insulator}.
	\newblock \emph{\bibinfo{journal}{Phys. Rev. B}} \textbf{\bibinfo{volume}{93}},
	\bibinfo{pages}{085438} (\bibinfo{year}{2016}).
	
	\bibitem{Kartashov-OL-2016}
	\bibinfo{author}{Kartashov, Y.~V.} \& \bibinfo{author}{Skryabin, D.~V.}
	\newblock \bibinfo{title}{Two-dimensional lattice solitons in polariton
		condensates with spin-orbit coupling}.
	\newblock \emph{\bibinfo{journal}{Opt. Lett.}} \textbf{\bibinfo{volume}{41}},
	\bibinfo{pages}{5043--5046} (\bibinfo{year}{2016}).
	
	\bibitem{DGulevich-Meissner}
	\bibinfo{author}{Gulevich, D.~R.}, \bibinfo{author}{Skryabin, D.~V.},
	\bibinfo{author}{Alodjants, A.~P.} \& \bibinfo{author}{Shelykh, I.~A.}
	\newblock \bibinfo{title}{Topological spin meissner effect in spinor
		exciton-polariton condensate: Constant amplitude solutions, half-vortices,
		and symmetry breaking}.
	\newblock \emph{\bibinfo{journal}{Phys. Rev. B}} \textbf{\bibinfo{volume}{94}},
	\bibinfo{pages}{115407} (\bibinfo{year}{2016}).
	
	\bibitem{Ablowitz-PRA-2014}
	\bibinfo{author}{Ablowitz, M.~J.}, \bibinfo{author}{Curtis, C.~W.} \&
	\bibinfo{author}{Ma, Y.-P.}
	\newblock \bibinfo{title}{Linear and nonlinear traveling edge waves in optical
		honeycomb lattices}.
	\newblock \emph{\bibinfo{journal}{Phys. Rev. A}} \textbf{\bibinfo{volume}{90}},
	\bibinfo{pages}{023813} (\bibinfo{year}{2014}).
	
	\bibitem{Leykam-PRL-2016}
	\bibinfo{author}{Leykam, D.} \& \bibinfo{author}{Chong, Y.~D.}
	\newblock \bibinfo{title}{Edge solitons in nonlinear-photonic topological
		insulators}.
	\newblock \emph{\bibinfo{journal}{Phys. Rev. Lett.}}
	\textbf{\bibinfo{volume}{117}}, \bibinfo{pages}{143901}
	(\bibinfo{year}{2016}).
	
	\bibitem{Kartashov-Optica-2016}
	\bibinfo{author}{Kartashov, Y.~V.} \& \bibinfo{author}{Skryabin, D.~V.}
	\newblock \bibinfo{title}{Modulational instability and solitary waves in
		polariton topological insulators}.
	\newblock \emph{\bibinfo{journal}{Optica 3}} \bibinfo{pages}{1228}
	(\bibinfo{year}{2016}).
	
	\bibitem{CedraMendez2016}
	\bibinfo{author}{Cerda-M\'endez, E.~A.} \emph{et~al.}
	\newblock \bibinfo{title}{Exciton-polariton gap solitons in two-dimensional
		lattices}.
	\newblock \emph{\bibinfo{journal}{Phys. Rev. Lett.}}
	\textbf{\bibinfo{volume}{111}}, \bibinfo{pages}{146401}
	(\bibinfo{year}{2013}).
	
	\bibitem{Soln-PRL-2017}
	\bibinfo{author}{Solnyshkov, D.~D.}, \bibinfo{author}{Bleu, O.},
	\bibinfo{author}{Teklu, B.} \& \bibinfo{author}{Malpuech, G.}
	\newblock \bibinfo{title}{Chirality of topological gap solitons in bosonic
		dimer chains}.
	\newblock \emph{\bibinfo{journal}{Phys. Rev. Lett.}}
	\textbf{\bibinfo{volume}{118}}, \bibinfo{pages}{023901}
	(\bibinfo{year}{2017}).
	
	\bibitem{Rajaraman}
	\bibinfo{author}{Rajaraman, R.}
	\newblock \emph{\bibinfo{title}{Solitons and Instantons: An Introduction to
			Solitons and Instantons in Quantum Field Theory}}
	(\bibinfo{publisher}{Elsevier Science}, \bibinfo{year}{1987}).
	
	\bibitem{Vasc-APL-2011}
	\bibinfo{author}{de~Vasconcellos, S.~M.} \emph{et~al.}
	\newblock \bibinfo{title}{Spatial, spectral, and polarization properties of
		coupled micropillar cavities}.
	\newblock \emph{\bibinfo{journal}{Appl. Phys. Lett.}}
	\textbf{\bibinfo{volume}{99}}, \bibinfo{pages}{101103}
	(\bibinfo{year}{2011}).
	
	\bibitem{Nalitov-PRL-2015}
	\bibinfo{author}{Nalitov, A.~V.}, \bibinfo{author}{Malpuech, G.},
	\bibinfo{author}{Ter\ifmmode~\mbox{\c{c}}\else \c{c}\fi{}as, H.} \&
	\bibinfo{author}{Solnyshkov, D.~D.}
	\newblock \bibinfo{title}{Spin-orbit coupling and the optical spin hall effect
		in photonic graphene}.
	\newblock \emph{\bibinfo{journal}{Phys. Rev. Lett.}}
	\textbf{\bibinfo{volume}{114}}, \bibinfo{pages}{026803}
	(\bibinfo{year}{2015}).
	
	\bibitem{Sala-PRX-2015}
	\bibinfo{author}{Sala, V.~G.} \emph{et~al.}
	\newblock \bibinfo{title}{Spin-orbit coupling for photons and polaritons in
		microstructures}.
	\newblock \emph{\bibinfo{journal}{Phys. Rev. X}} \textbf{\bibinfo{volume}{5}},
	\bibinfo{pages}{011034} (\bibinfo{year}{2015}).
	
	\bibitem{Hasan-RMP-2010}
	\bibinfo{author}{Hasan, M.~Z.} \& \bibinfo{author}{Kane, C.~L.}
	\newblock \bibinfo{title}{Colloquium: Topological insulators}.
	\newblock \emph{\bibinfo{journal}{Rev. Mod. Phys.}}
	\textbf{\bibinfo{volume}{82}}, \bibinfo{pages}{3045--3067}
	(\bibinfo{year}{2010}).
	
	\bibitem{Fruchart}
	\bibinfo{author}{Fruchart, M.} \& \bibinfo{author}{Carpentier, D.}
	\newblock \bibinfo{title}{An introduction to topological insulators}.
	\newblock \emph{\bibinfo{journal}{Comptes Rendus Physique}}
	\textbf{\bibinfo{volume}{14}}, \bibinfo{pages}{779 -- 815}
	(\bibinfo{year}{2013}).
	\newblock \bibinfo{note}{Topological insulators / Isolants
		topologiquesTopological insulators / Isolants topologiques}.
	
	\bibitem{Vladimirova}
	\bibinfo{author}{Vladimirova, M.} \emph{et~al.}
	\newblock \bibinfo{title}{Polarization controlled nonlinear transmission of
		light through semiconductor microcavities}.
	\newblock \emph{\bibinfo{journal}{Phys. Rev. B}} \textbf{\bibinfo{volume}{79}},
	\bibinfo{pages}{115325} (\bibinfo{year}{2009}).
	
	\bibitem{LL-III}
	\bibinfo{author}{Landau, L.~D.} \& \bibinfo{author}{Lifshitz, L.~M.}
	\newblock \emph{\bibinfo{title}{Quantum Mechanics, Third Edition:
			Non-Relativistic Theory (Volume 3)}} (\bibinfo{publisher}{Pregamon Press},
	\bibinfo{year}{1977}).
	
	\bibitem{Ablowitz}
	\bibinfo{author}{Ablowitz, M.~J.}
	\newblock \emph{\bibinfo{title}{Nonlinear Dispersive Waves: Asymptotic Analysis
			and Solitons}} (\bibinfo{publisher}{Cambridge University Press},
	\bibinfo{year}{2011}).
	
	\bibitem{Milicevic}
	\bibinfo{author}{Milićević, M.} \emph{et~al.}
	\newblock \bibinfo{title}{Edge states in polariton honeycomb lattices}.
	\newblock \emph{\bibinfo{journal}{2D Materials}} \textbf{\bibinfo{volume}{2}},
	\bibinfo{pages}{034012} (\bibinfo{year}{2015}).
	
	\bibitem{Galbiati}
	\bibinfo{author}{Galbiati, M.} \emph{et~al.}
	\newblock \bibinfo{title}{Polariton condensation in photonic molecules}.
	\newblock \emph{\bibinfo{journal}{Phys. Rev. Lett.}}
	\textbf{\bibinfo{volume}{108}}, \bibinfo{pages}{126403}
	(\bibinfo{year}{2012}).
	
	\bibitem{Abbarchi}
	\bibinfo{author}{Abbarchi, M.} \emph{et~al.}
	\newblock \bibinfo{title}{Macroscopic quantum self-trapping and josephson
		oscillations of exciton polaritons}.
	\newblock \emph{\bibinfo{journal}{Nature Physics}}
	\textbf{\bibinfo{volume}{9}}, \bibinfo{pages}{275–279}
	(\bibinfo{year}{2013}).
	
	\bibitem{Duff-APL}
	\bibinfo{author}{Dufferwiel, S.} \emph{et~al.}
	\newblock \bibinfo{title}{Tunable polaritonic molecules in an open microcavity
		system}.
	\newblock \emph{\bibinfo{journal}{Applied Physics Letters}}
	\textbf{\bibinfo{volume}{107}}, \bibinfo{pages}{201106}
	(\bibinfo{year}{2015}).
	
	\bibitem{Kuwata-PRL-1997}
	\bibinfo{author}{Kuwata-Gonokami, M.} \emph{et~al.}
	\newblock \bibinfo{title}{Parametric scattering of cavity polaritons}.
	\newblock \emph{\bibinfo{journal}{Phys. Rev. Lett.}}
	\textbf{\bibinfo{volume}{79}}, \bibinfo{pages}{1341--1344}
	(\bibinfo{year}{1997}).
	
	\bibitem{Savvidis-PRL-2000}
	\bibinfo{author}{Savvidis, P.~G.} \emph{et~al.}
	\newblock \bibinfo{title}{Angle-resonant stimulated polariton amplifier}.
	\newblock \emph{\bibinfo{journal}{Phys. Rev. Lett.}}
	\textbf{\bibinfo{volume}{84}}, \bibinfo{pages}{1547--1550}
	(\bibinfo{year}{2000}).
	
	\bibitem{Ciuti-PRB-2001}
	\bibinfo{author}{Ciuti, C.}, \bibinfo{author}{Schwendimann, P.} \&
	\bibinfo{author}{Quattropani, A.}
	\newblock \bibinfo{title}{Parametric luminescence of microcavity polaritons}.
	\newblock \emph{\bibinfo{journal}{Phys. Rev. B}} \textbf{\bibinfo{volume}{63}},
	\bibinfo{pages}{041303} (\bibinfo{year}{2001}).
	
	\bibitem{Dasbach-PRB-2001}
	\bibinfo{author}{Dasbach, G.}, \bibinfo{author}{Schwab, M.},
	\bibinfo{author}{Bayer, M.} \& \bibinfo{author}{Forchel, A.}
	\newblock \bibinfo{title}{Parametric polariton scattering in microresonators
		with three-dimensional optical confinement}.
	\newblock \emph{\bibinfo{journal}{Phys. Rev. B}} \textbf{\bibinfo{volume}{64}},
	\bibinfo{pages}{201309} (\bibinfo{year}{2001}).
	
	\bibitem{Tartakovskii-PRB-2002}
	\bibinfo{author}{Tartakovskii, A.~I.} \emph{et~al.}
	\newblock \bibinfo{title}{Polariton parametric scattering processes in
		semiconductor microcavities observed in continuous wave experiments}.
	\newblock \emph{\bibinfo{journal}{Phys. Rev. B}} \textbf{\bibinfo{volume}{65}},
	\bibinfo{pages}{081308} (\bibinfo{year}{2002}).
	
\end{thebibliography}

\end{document}
